\subsection{Front-end}
\label{subsec:front-end}

\subsubsection{Chart.js}
\label{subsubsec:chart}
Chart.js è una libreria JavaScript che permette di integrare diverse tipologie di grafici (grafici a torta, grafici a barre, grafici di dispersione) all'interno di una pagina web tramite l'utilizzo dei HTML5 canvans.\\
\textbf{Motivazioni}
\begin{enumerate}
    \item \textbf{Popolarità}: è la libreria di grafici più utilizzata;
    \item \textbf{Integrazione}: è compatibile con Angular per cui esiste un apposito pacchetto che fornisce delle direttive per ogni tipologia di grafico.
\end{enumerate}